\documentclass[12pt]{article}
\usepackage{amsmath,amssymb,amsfonts}
\usepackage[margin=1in]{geometry}

\begin{document}

\begin{center}
{\Large\bfseries Chapter 10, Problem 2}
\end{center}

\noindent\textbf{Problem.} Consider the group built up from the unit matrix and the Pauli spin matrices. The elements of this group are $\pm I$, $\pm i\sigma_x$, $\pm i\sigma_y$, and $\pm i\sigma_z$, where
\[
\sigma_x = \begin{pmatrix} 0 & 1 \\ 1 & 0 \end{pmatrix}, \quad
\sigma_y = \begin{pmatrix} 0 & -i \\ i & 0 \end{pmatrix}, \quad
\sigma_z = \begin{pmatrix} 1 & 0 \\ 0 & -1 \end{pmatrix}.
\]
What is the multiplication table? Find all subgroups and classes. Are there any invariant subgroups?

\bigskip
\noindent\textbf{Solution.}

The eight elements of the group are:
\[
G = \{I, -I, i\sigma_x, -i\sigma_x, i\sigma_y, -i\sigma_y, i\sigma_z, -i\sigma_z\}.
\]
For brevity, denote these as $\{e, \bar{e}, a, \bar{a}, b, \bar{b}, c, \bar{c}\}$ where $\bar{g} = -g$.

The key multiplication rules come from the Pauli matrix algebra:
\[
\sigma_x \sigma_y = i\sigma_z, \quad \sigma_y \sigma_z = i\sigma_x, \quad \sigma_z \sigma_x = i\sigma_y,
\]
and $\sigma_j^2 = I$ for $j = x, y, z$. Therefore:
\[
(i\sigma_x)^2 = -I, \quad (i\sigma_y)^2 = -I, \quad (i\sigma_z)^2 = -I,
\]
\[
(i\sigma_x)(i\sigma_y) = -\sigma_x\sigma_y = -i\sigma_z, \quad
(i\sigma_y)(i\sigma_x) = -\sigma_y\sigma_x = i\sigma_z.
\]

\noindent\textbf{Multiplication Table} (using $e, \bar{e}, a, \bar{a}, b, \bar{b}, c, \bar{c}$):

\[
\begin{array}{c|cccccccc}
 & e & \bar{e} & a & \bar{a} & b & \bar{b} & c & \bar{c} \\
\hline
e & e & \bar{e} & a & \bar{a} & b & \bar{b} & c & \bar{c} \\
\bar{e} & \bar{e} & e & \bar{a} & a & \bar{b} & b & \bar{c} & c \\
a & a & \bar{a} & \bar{e} & e & \bar{c} & c & b & \bar{b} \\
\bar{a} & \bar{a} & a & e & \bar{e} & c & \bar{c} & \bar{b} & b \\
b & b & \bar{b} & c & \bar{c} & \bar{e} & e & \bar{a} & a \\
\bar{b} & \bar{b} & b & \bar{c} & c & e & \bar{e} & a & \bar{a} \\
c & c & \bar{c} & \bar{b} & b & a & \bar{a} & \bar{e} & e \\
\bar{c} & \bar{c} & c & b & \bar{b} & \bar{a} & a & e & \bar{e}
\end{array}
\]

\medskip
\noindent\textbf{Subgroups:}

By Lagrange's theorem, proper subgroups of a group of order 8 must have order 2 or 4.

\textit{Order 2:} $\{e, \bar{e}\} = \{I, -I\}$. This is the only subgroup of order 2 (since $\bar{e}$ is the only element of order 2).

\textit{Order 4:}
\begin{align*}
H_1 &= \{e, \bar{e}, a, \bar{a}\} = \{I, -I, i\sigma_x, -i\sigma_x\}, \\
H_2 &= \{e, \bar{e}, b, \bar{b}\} = \{I, -I, i\sigma_y, -i\sigma_y\}, \\
H_3 &= \{e, \bar{e}, c, \bar{c}\} = \{I, -I, i\sigma_z, -i\sigma_z\}.
\end{align*}
Each is isomorphic to the cyclic group $\mathbb{Z}_4$.

\medskip
\noindent\textbf{Conjugacy Classes:}

Since $-I$ commutes with everything, $\{-I\}$ is its own class. For the others, we compute, e.g.:
\[
(i\sigma_y)(i\sigma_x)(i\sigma_y)^{-1} = (i\sigma_y)(i\sigma_x)(-i\sigma_y) = -i\sigma_x.
\]
So $i\sigma_x$ and $-i\sigma_x$ are conjugate. Similarly for the others.

The five conjugacy classes are:
\[
\{e\}, \quad \{\bar{e}\}, \quad \{a, \bar{a}\}, \quad \{b, \bar{b}\}, \quad \{c, \bar{c}\}.
\]

\medskip
\noindent\textbf{Invariant Subgroups:}

A subgroup is invariant (normal) if it is a union of complete conjugacy classes. Since the classes are $\{e\}$, $\{\bar{e}\}$, $\{a,\bar{a}\}$, $\{b,\bar{b}\}$, $\{c,\bar{c}\}$:

\begin{itemize}
\item $\{e, \bar{e}\}$ is invariant (union of $\{e\}$ and $\{\bar{e}\}$, order 2 divides 8). \checkmark
\item $H_1 = \{e, \bar{e}, a, \bar{a}\}$ is invariant (union of $\{e\}$, $\{\bar{e}\}$, $\{a,\bar{a}\}$). \checkmark
\item $H_2 = \{e, \bar{e}, b, \bar{b}\}$ is invariant. \checkmark
\item $H_3 = \{e, \bar{e}, c, \bar{c}\}$ is invariant. \checkmark
\end{itemize}

All four proper nontrivial subgroups are invariant. This group is known as the \textbf{quaternion group} $Q_8$.

\end{document}
